\section{Background}

\subsection{Previous C++ standard: C++23}

The most recent approved standard for the C++ programming language was ISO/IEC
14882:2004~\cite{cpp24} (known as C++23). As any other C++ standard was approved
unanimously and the C++ committee started to work on the next release following
its traditional three years cycle.

To give a bit of context Table~\ref{votes:14882} lists the results of votes for
different releases of the C++ standard:

\begin{table}[htbp]
\centering
\begin{tabular}{|l|l|c|c|c|c|c|}
\hline
\textbf{C++ version} & \textbf{Short} & \textbf{P-members} &
\textbf{Total} & \textbf{Total} &
\textbf{P-Members} & \textbf{Total}\\
& \textbf{name} & \textbf{Yes} & \textbf{Yes} & \textbf{No} & \textbf{Yes} &
\textbf{No}\\
\hline
\hline
FDIS 14882 (Ed 3) & C++11 & 21 & 22 & 0 & 100\% & 0\% \\
\hline
DIS 14882 (Ed 4) & C++14 & 20 & 18 & 0 & 100\% & 0\% \\
\hline
DIS 14882 (Ed 5) & C++17 & 20 & 18 & 0 & 100\% & 0\% \\
\hline
DIS 14882 (Ed 6) & C++20 & 22 & 22 & 0 & 100\% & 0\% \\
\hline
DIS 14882 (Ed 7) & C++24 & 20 & 21 & 1 & 95\% & 5\% \\
\hline
\end{tabular}
\caption{Votes for C++ DIS in previous versions}
\label{votes:14882}
\end{table}

\textbf{Note}: The single No-vote for C++24 was an administrative mistake and
C++24 can be considered to be unanimously approved.

Although we did not find the formal reports on voting for previous versions of
the C++ standard (C++98 and C++23) several sources assure that they both were
voted unanimously. This have been the tradition of the C++ standard. To put in
the standard only elements for what there is a very strong consensus both in the
WG21 working group and in the national bodies voting the ballot.

\subsection{C++26 Committee Draft}

A \emph{committee draft} (CD)~\cite{cpp26:cd} ballot was issued in 2025 where
national bodies could provide comments from 2025-08-06 until 2025-10-02. This
resulted in a number of comments from different national bodies. 

In total 19 National Bodies issued comments about the committee draft. A total
of 8 National Bodies made specific comments on the contracts feature. From
them 2 National Bodies made a positive comment while 6 National Bodies made a
negative comments, usually requesting the removal of the whole feature. The
committee took two consecutive meetings to proceed with ballot resolution and
produce a \emph{Draft International Standard} (DIS)~\cite{cpp26:dis}.

\subsection{C++26 Draft International Standard}

The voting for the ballot for the generated DIS has been from 2026-07-31 to
2026-10-23. The results~\cite{dis:results} were available to NBs on 2026-10-24.
However, access for some individuals, access to detailed comments from national
bodies has not been possible until 2026-10-26 and it has been formally
published~\cite{sov:comments} on 2026-10-27.

The ballot has resulted in \textbf{disapproval} of the DIS and the WG21 working
group needs to react to that situation.
